\documentclass[11pt]{article}
% TrustBridge paper draft v0.1 (2026-09-28)
% Single-column draft for internal review; final version will be re-typeset
% in the venue template (LNCS for MICCAI workshop / IEEEtran for ISBI).
\usepackage[margin=1in]{geometry}
\usepackage{lmodern}
\usepackage{amssymb}
\usepackage{graphicx}
\usepackage{booktabs}
\usepackage{multirow}
\usepackage{amsmath}
\usepackage[numbers,sort&compress]{natbib}
\usepackage[colorlinks=true,linkcolor=blue,citecolor=blue,urlcolor=blue]{hyperref}
\usepackage{xcolor}

% ---- terminology placeholder: rename RCRF globally once decided ----
\newcommand{\rcrf}{recursive convergence residual field (RCRF)}
\newcommand{\TrustMap}{trust map}
\newcommand{\TrustScore}{Trust Score}

\title{TrustBridge: Trust-Map Quantification for Diffusion-Bridge\\
Medical Image Translation via Recursive Convergence Residuals}

\author{%
  Author list to be finalized\\
  Institution\\
  \small Code and weights: \url{https://github.com/ASHDUSK/TrustBridge-}
}
\date{September 2026}

\begin{document}
\maketitle

\begin{abstract}
Diffusion bridges achieve state-of-the-art medical image translation, yet their
outputs carry no indication of \emph{where} the synthesis can be trusted---a gap
that the original SelfRDB authors explicitly identify as an open problem.
We present TrustBridge, a training-free trust-quantification layer that harvests
two signals already present inside the self-consistent recursive bridge sampler:
(1) the \emph{recursive convergence residual field} (RCRF), the per-pixel
discrepancy between successive recursive target estimates, marking pixels the
model repeatedly revisits without reaching self-consistency; and (2)
\emph{stochastic-bridge ensemble uncertainty}, the per-pixel dispersion across
independent bridge chains arising from the stochastic endpoints and latent
variables native to the sampler. The two signals are fused into a pixel-wise
\TrustMap{}, aggregated into an image-level \TrustScore{} (0--100), and gated by
a three-tier rule (PASS/REVIEW/REJECT) with an input-modality check. On held-out
data spanning three datasets and five translation directions, the \TrustMap{}
correlates strongly with true pixel error (Spearman $-0.81$ to $-0.90$);
low-trust pixels localize high-error pixels with AUROC up to $0.93$; the
least-trusted error decile exceeds the most-trusted decile by $35$--$73\times$;
and discarding the least-trusted 20\% of pixels improves PSNR by up to
$6.95$\,dB on real tumor patients. Ablations confirm that each module is
individually effective and their fusion is optimal. Five controlled experiments
(adaptive resampling, learnable aggregation, image-level regression, intensity
harmonization, and round-trip consistency) further delineate the boundary of
the approach. The system ships as an open-source clinical workbench with
fully reproducible training pipelines.
\end{abstract}

\section{Introduction}
\label{sec:intro}
Multi-contrast magnetic resonance imaging (MRI) captures complementary tissue
contrast, but acquiring every clinically relevant weighting is expensive, slow,
and sometimes contraindicated. Medical image translation---synthesizing an
unacquired modality from an acquired one---addresses this scarcity and now
achieves excellent fidelity. Denoising diffusion probabilistic models
(DDPMs)~\cite{ho2020ddpm,nichol2021iddpm} replaced GANs in this task for their
training stability, and diffusion bridges connect the source and target
distributions directly rather than through an isotropic-noise bottleneck. In
particular, SelfRDB~\cite{arslan2025selfrdb} introduced a self-consistent
recursive diffusion bridge that reaches state-of-the-art multi-contrast
translation.

Trust, however, has not kept pace with fidelity. Because translation is
ill-posed, generative models can \emph{hallucinate}: synthesize plausible
structure that is absent from the underlying anatomy. In clinical use, a
radiologist cannot distinguish such fabrication from faithful synthesis, and
this trust bottleneck---not image quality---is the primary obstacle to
deployment. The SelfRDB authors state this gap explicitly: the uncertainty of
diffusion bridges ``is warranted to [be] evaluate[d]'' in future work
\cite{arslan2025selfrdb}.

We close this gap with an observation: the self-consistent recursive sampler
already computes, at every step and every pixel, the evidence needed to assess
its own reliability---it simply discards it. SelfRDB reduces the discrepancy
between successive recursive estimates to a single stopping scalar; the
per-pixel residual field carries the spatial distribution of the model's
\emph{inability to self-consist}. Likewise, the bridge endpoint noise and
generator latent variables make repeated runs differ; their dispersion marks
where the model holds no stable belief. TrustBridge systematizes both signals
into a \TrustMap{}, an image-level \TrustScore{}, and an archivable QC report,
with three properties that distinguish it from prior uncertainty machinery:
(i) \emph{training-free}---no model change, no labels, no extra network;
(ii) \emph{sampler-native}---signals are extracted from quantities the sampler
already computes; and (iii) \emph{pixel-level}---unlike image-level hallucination
metrics~\cite{tivnan2024hallucinationindex}, the output supports a radiologist's
local reading.

\paragraph{Contributions.}
\begin{enumerate}
  \item We propose the \rcrf{} (RCRF) and stochastic-bridge ensemble
        uncertainty as complementary, training-free trust signals for
        diffusion-bridge translation, fused into a calibrated \TrustMap{} with
        a three-tier QC gate and an input-modality gate
        (\S\ref{sec:method}).
  \item We validate the \TrustMap{} on three public datasets and five
        translation directions: error localization AUROC up to $0.93$,
        a $35$--$73\times$ error separation between the least- and
        most-trusted deciles, and sparsification gains up to $+6.95$\,dB on
        \emph{real tumor patients} with co-registered ground truth
        (\S\ref{sec:exp}).
  \item We report five controlled experiments whose outcomes are mixed or
        negative---adaptive resampling, learnable aggregation, image-level
        regression, intensity harmonization, and round-trip consistency---and
        use them to delineate the method's boundary: at this data scale, the
        pixel-wise \TrustMap{} is the only reliable product form, and
        learning-based aggregation overfits six-subject validation sets
        (\S\ref{sec:exp}, \S\ref{sec:discussion}).
\end{enumerate}

\section{Related Work}
\label{sec:related}
\section{Related Work}
\label{sec:related}

\subsection{Medical Image Translation}
Unpaired cross-modality synthesis has evolved from adversarial translation
---CycleGAN~\cite{zhu2017cyclegan} and pix2pix~\cite{isola2017pix2pix} established
the unpaired and paired regimes, followed by multi-contrast MRI specialists such
as conditional GANs~\cite{dar2019cgan} and hybrid CNN--transformers
(ResViT~\cite{dalmaz2022resvit})---to diffusion-based models. SynDiff
\cite{ozbey2023syndiff} introduced adversarial diffusion with a cyclic
consistency prior for unsupervised medical translation, and
I2I-Galip~\cite{korkmaz2025i2igalip} replaced the discriminator with a
vision--language model for efficiency. In parallel, bridges connect the two
image distributions directly rather than through an isotropic-noise bottleneck:
DDPMs~\cite{ho2020ddpm} and their score-based generalization~\cite{song2021scoresde}
provide the mathematical backbone, with score-based
NCSN~\cite{song2019ncsn}, improved schedules~\cite{nichol2021iddpm}, direct
diffusion bridges with data consistency~\cite{chung2023directbridge}, dual
implicit bridges~\cite{su2023ddib}, neural Schr{\"o}dinger
bridges~\cite{kim2024nsb}, and direct iterative inversion~\cite{delbracio2023indi}
spanning the design space. SelfRDB~\cite{arslan2025selfrdb} combines a soft-prior
bridge with self-consistent recursive sampling and reports state-of-the-art
multi-contrast translation, while explicitly leaving \emph{uncertainty
quantification} of the bridge as an open problem. Our work builds directly on
this sampler: we keep the model frozen and harvest two signals that SelfRDB
discards.

\subsection{Uncertainty Quantification in Medical Imaging}
Per-pixel uncertainty is commonly obtained by MC-Dropout~\cite{gal2016mcdropout},
deep ensembles~\cite{lakshminarayanan2017ensembles}, or test-time
augmentation~\cite{ayhan2018tta}; the epistemic/aleatoric taxonomy of
\cite{kendall2017uncertainties} frames these as sampling-based epistemic
estimates. Diffusion ensembles follow the same recipe: implicit segmentation
ensembles~\cite{wolleb2022diffusionensembles} and per-pixel standard deviations
across diffusion restorations~\cite{osuna2024microscopy} established that
inter-sample dispersion tracks difficult image regions. We inherit this insight
but differ in three ways: (i) our ensemble members are sampled from a
\emph{diffusion bridge} whose stochastic endpoints and latent variables are
native to the sampler, not injected ad hoc; (ii) we pair ensemble dispersion
with a second, complementary signal---the convergence residual---that
is unique to self-consistent recursive samplers and costs nothing; and (iii) we
fuse both into a calibrated trust map with a clinical QC gate rather than
reporting dispersion alone. Notably, concurrent bridge models actively
\emph{suppress} this stochasticity for reconstruction
\cite{song2024ccbm}; we argue it is an asset when the task is trust assessment,
even though it is a defect for point synthesis.

\subsection{Hallucination Assessment and Trustworthy Radiology AI}
Closest to our goal, the Hallucination Index~\cite{tivnan2024hallucinationindex}
quantifies fabricated content of generative reconstruction models as an
\emph{image-level} metric over repeated samples; MUMINS~\cite{oliveras2026mumins}
conditions medical image synthesis on metadata and trains an uncertainty-aware
prediction head to output a spatial uncertainty map; and the measurement-geometry
analysis of \cite{jin2026trustworthy} proves when a measurement operator can
distinguish solutions, explaining hallucination from an operator-design
perspective. These approaches either require repeated sampling at image
granularity, architectural modification with retraining, or ex-ante operator
design. TrustBridge instead derives pixel-level evidence \emph{inside} the
existing sampler, requires no architectural change or labels, and additionally
supports rule-based gating: our framework is complementary to output-guardrail
lifecycles such as RAISE~\cite{cardoso2023raise} and provides the objective,
pixel-level evidence that reporting frameworks like
AI-RADS~\cite{fink2026airads} call for, while remaining fully automatic.

\paragraph{Beyond sampling dispersion: semantic and pathology-aware signals.}
Dispersion-style evidence is blind to errors the sampler makes
\emph{consistently}---a limitation we quantify in \S\ref{sec:exp}. Three
concurrent research lines attack this from the evaluation side and mark the
natural next step for trust layers like ours. CRAFT~\cite{chung2026craft}
scores generated images against clinical criteria with a foundation-model
Clinical Alignment Score, reducing hallucination-like content; ViCTr~\cite{susladkar2025victr}
makes synthesis pathology-aware with graded severity control via EWC-regularized
LoRA adaptation; and DDWM~\cite{li2025ddwm} achieves zero-shot translation by
learning only the target-domain distribution and fusing domain-invariant
anatomy through wavelet decomposition. All three require either foundation
models, retraining, or architectural change; they are complementary to
sampler-native evidence and are discussed as future work in \S\ref{sec:discussion}.


\section{Method}
\label{sec:method}

\subsection{Preliminaries: the self-consistent recursive diffusion bridge}
SelfRDB~\cite{arslan2025selfrdb} constructs a bridge whose forward process
starts at the target image $x_0$ and ends at a noise-added source image
$y_\epsilon$, with intermediate samples
$x_t=\mu_{x_0,t}x_0+\mu_{y,t}y+\sigma_t\epsilon$ where
$\sigma_t^2=\gamma s_t^2/(\bar s_t^2+s_t^2)$ grows monotonically toward the
source end (\emph{soft prior}). Reverse sampling starts at $x_T=y_\epsilon$;
at every step $t$, a generator
$G_\theta(x_t,t,y,\tilde x_0^{(r)})$ recursively refines a target estimate
$\tilde x_0^{(r+1)}$ until the change between successive estimates falls below
a threshold $\tau$ (SelfRDB: a $1\%$ global mean), after which the
self-consistent estimate $\tilde x_0^{\ast}$ drives a posterior draw of
$x_{t-1}$. The publicly released configuration uses $T{=}10$ bridge steps and
at most $R{=}2$ recursions, which makes per-slice inference interactive.

\subsection{Recursive convergence residual field (RCRF)}
Within one reverse step, the recursion produces a sequence of pixel-wise
discrepancies
$R^{(r)}(p)=\lvert\tilde x_0^{(r+1)}(p)-\tilde x_0^{(r)}(p)\rvert$, $p$ indexing
pixels. SelfRDB aggregates $R^{(r)}$ to a global mean for its stopping rule and
discards the spatial information. We retain the residual field over the final
$K$ reverse steps (the steps closest to the data distribution),
\begin{equation}
  \mathrm{RCRF}(p) \;=\; \frac{1}{K}\sum_{t\in\mathcal{K}}\;\sum_{r}
      \bigl\lvert \tilde x_0^{(r+1)}(p)-\tilde x_0^{(r)}(p)\bigr\rvert ,
  \label{eq:rcrf}
\end{equation}
where $\mathcal{K}$ is the capture window ($K{=}3$ here). Intuition: pixels the
generator keeps revisiting without converging are pixels for which its learned
prior provides insufficient evidence---prime candidates for hallucinated
content. The signal is epistemic in the sense
of~\cite{kendall2017uncertainties} and costs nothing: the recursion runs
anyway.

\subsection{Stochastic-bridge ensemble uncertainty}
The bridge endpoint $y_\epsilon$ and the generator's internal latent variables
$z$ make two runs of the sampler on the same input differ. Running
$N$ independent chains (distinct endpoint noise and latents) and taking the
per-pixel standard deviation,
$U(p)=\operatorname{std}_{n}\{x_0^{(n)}(p)\}$, yields a second signal that
measures \emph{sampling consistency}: pixels whose synthetic value is unstable
across chains are regions where the model holds no stable belief. Unlike
MC-Dropout~\cite{gal2016mcdropout} or deep
ensembles~\cite{lakshminarayanan2017ensembles}, no dropout is injected and no
separate networks are trained; the dispersion arises from the bridge sampler
itself. Prior diffusion-ensemble work in segmentation and
restoration~\cite{wolleb2022diffusionensembles,osuna2024microscopy} established
that such dispersion tracks difficult regions; we pair it with a
non-dispersion signal (RCRF) that is orthogonal to sampling variance.

\subsection{Trust map, Trust Score, and QC gate}
Both fields are robustly normalized (99th-percentile clipping) and fused as
$\TrustMap(p)=1-w_1\widehat{\mathrm{RCRF}}(p)-w_2\widehat{U}(p)$ with
$w_1{=}0.6$, $w_2{=}0.4$ (sensitivity analyzed in \S\ref{sec:exp}). The image
level $\TrustScore$ is the body-region mean of $\TrustMap$ scaled to 0--100,
gated at $75$ (PASS), $55$ (REVIEW), and $55$ (REJECT). A per-slice report
(PNG + JSON) records the maps, the per-step recursion statistics, and the
input-gate verdict. An input-modality gate---a light classifier on
body-standardized percentile and gradient features---flags inputs whose
characteristics do not match the declared source modality before any synthesis
occurs; on IXI test data it verifies correctly for $99$--$100\%$ of matched
inputs and detects $99$--$100\%$ of mismatched ones.

\section{Experiments}
\label{sec:exp}

\subsection{Setup}
\paragraph{Datasets.}
\textbf{IXI}~\cite{ixi_dataset}: 40 healthy subjects (T1/T2/PD), subject-wise
split 26/6/8 with no subject overlap; $4{,}000{+}$ paired axial slices under
the official protocol (per-slice statistics, official mean-norm PSNR/SSIM).
\textbf{BraTS2018}~\cite{menze2015brats}: 6 HGG patients (399 tumor-bearing
slices) with real co-registered T2 ground truth and tumor segmentation; source
FLAIR. \textbf{SynthRAD2023 pelvis}~\cite{thummerer2023synthrad}: 10 patients,
center-matched T1w$\to$CT and T2w$\to$CT.
\paragraph{Models.}
Two self-trained 128$\times$128 bridges (T2$\to$T1 and T1$\to$T2, 30 epochs,
single RTX 4060 Laptop 8\,GB, recipe matched to the official release) provide
the quantitative backbone; official SelfRDB weights provide 256$\times$256
demonstration and the BraTS FLAIR$\to$T2 evaluation.
\paragraph{Ethics.}
All three datasets are public research datasets; all processing is local; no
output is used for clinical decisions.

\subsection{Translation quality and trust localization}
\begin{table}[t]
\centering\small
\caption{Main results. IXI: self-trained 128$\times$128 models, official
mean-norm PSNR/SSIM protocol. ``Corrected'' removes a global linear intensity
mismatch (see \S\ref{sec:discussion}). BraTS: official
FLAIR$\to$T2 weights, 399 tumor-bearing slices.}
\label{tab:main}
\begin{tabular}{llcccc}
\toprule
Dataset & Direction & PSNR (dB) & SSIM (\%) & NCC & Corrected PSNR \\
\midrule
IXI (paper value~\cite{arslan2025selfrdb}) & T2$\to$T1 & 31.56 & 95.65 & -- & -- \\
IXI, ours & T2$\to$T1 & 19.81$\pm$1.10 & 71.0$\pm$6.6 & 0.884 & 31.74 \\
IXI, ours & T1$\to$T2 & 20.44$\pm$1.16 & 80.4$\pm$4.1 & 0.905 & 32.52 \\
BraTS2018, official & FLAIR$\to$T2 & 24.06$\pm$1.62 & 87.38 & 0.944 & -- \\
\bottomrule
\end{tabular}
\end{table}

Table~\ref{tab:main} reports translation quality. On tumor patients, the
official FLAIR$\to$T2 weights reach $23.9$\,dB against real co-registered T2
ground truth---the first successful quantitative validation of bridge
translation on \emph{pathological} data in our pipeline. On IXI, decomposing
the error shows that the content fidelity of our 128$\times$128 self-trained
models already matches the paper's setting: after removing a global linear
intensity mismatch, PSNR reaches $31.74$/$32.52$\,dB (T2$\to$T1/T1$\to$T2)
versus the paper's $31.56$\,dB, while the raw gap is attributable to a global
intensity offset and the reduced training budget.

\paragraph{Trust localization.}
Across IXI test slices, the \TrustMap{} correlates with true pixel error at
Spearman $\rho=-0.806$ (T2$\to$T1) and $\rho=-0.903$ (T1$\to$T2), and low-trust
pixels localize the top-error decile with AUROC $0.819$ and $0.902$
respectively. On BraTS tumor patients, trust AUROC is $0.932$.

\subsection{Ablation and controls}
\begin{table}[t]
\centering\small
\caption{Ablation and controls (7-row protocol). AUROC: low-trust pixels
predicting the top-error decile. $\Delta$PSNR: gain after dropping the
least-trusted 20\% of pixels.}
\label{tab:ablation}
\begin{tabular}{lcccc}
\toprule
Scheme & \multicolumn{2}{c}{AUROC} & \multicolumn{2}{c}{$\Delta$PSNR@20\% (dB)} \\
 & T2$\to$T1 & T1$\to$T2 & T2$\to$T1 & T1$\to$T2 \\
\midrule
Random (control) & 0.501 & 0.498 & $\approx$0 & $\approx$0 \\
Gradient baseline & 0.448 & 0.777 & +0.02 & +1.46 \\
Input-difference$^{\dagger}$ & -- & -- & +2.02 & +0.39 \\
RCRF only (N=1) & 0.792 & 0.880 & +2.34 & +4.56 \\
Ensemble only & 0.806 & 0.897 & +2.45 & +4.90 \\
\textbf{Fused (ours)} & \textbf{0.819} & \textbf{0.902} & \textbf{+2.14} & \textbf{+5.19} \\
Oracle (error map) & 1.000 & 1.000 & +10.13 & +15.03 \\
\bottomrule
\end{tabular}
\end{table}

Table~\ref{tab:ablation}: every module is individually effective far above the
random control; the fusion matches or beats each single signal; and the oracle
row quantifies the headroom---our signals capture $21$--$35\%$ of the
theoretically available sparsification gain. Weight sensitivity is mild
($<0.01$ AUROC across $w_1\in[0.3,0.6]$).

\subsection{Calibration}
Binning test pixels by trust decile yields a monotone reliability curve: the
least-trusted decile carries $35.3\times$ (T2$\to$T1) and $73.2\times$
(T1$\to$T2) the mean absolute error of the most-trusted decile, with $88.9\%$
of adjacent-bin transitions monotone. The trust map is therefore not merely a
ranking but a strongly discriminative magnitude estimate.

\subsection{Boundary experiments}
Five controlled experiments delineate what the layer can and cannot do at this
data scale:
(1)~\emph{learnable aggregation}: a 16-feature logistic/GBDT aggregator trained
on the validation split improves T2$\to$T1 marginally ($+0.010$ AUROC) but
regresses T1$\to$T2 ($-0.03$), overfitting the 6-subject validation set;
(2)~\emph{image-level regression}: a 16-feature ridge model predicts image
PSNR worse than the baseline trust mean ($0.32$--$0.33$ vs.\ $0.38$--$0.40$);
(3)~\emph{adaptive resampling} and (4)~\emph{round-trip consistency
signal}: no quality gain ($-0.016$\,dB and AUROC $0.503$ respectively) at
$35$--$75\%$ extra cost---consistent with ensemble-mean saturation;
(5)~\emph{intensity harmonization}: post-hoc quantile/affine matching degrades
raw PSNR ($-0.8$/$-0.9$\,dB) because the per-slice intensity offset is not
predictable without ground truth. These results are reported in full as
negative findings: at 40 subjects, the pixel-wise \TrustMap{} is the only
reliable product form, and learning-based aggregation requires
$\ge$20 validation subjects.

\subsection{Tumor-region analysis}
On the BraTS cohort, per-slice analysis shows that (i) pixel error inside the
tumor is $1.54\times$ the error outside, confirming that hallucination risk
concentrates at the lesion; and (ii) the \TrustMap{}---computed without any
tumor knowledge---separates tumor from non-tumor pixels with per-slice AUROC
$0.871$, capturing $94\%$ of the localization quality of the error map itself
($0.925$). This supports the clinical reading: flagged regions coincide with
the places a radiologist must scrutinize.

\section{Discussion}
\label{sec:discussion}
\paragraph{Reproducibility.}
Officially released IXI weights reproduce only $\sim$15\,dB on independently
prepared IXI data (the paper reports 31.56\,dB). A systematic exclusion study
(rigid misalignment, $z$-shift, intensity window, cropping, orientation)
traces the gap to preprocessing details that the original authors do not
publish; the same issue is independently reported for the companion SynDiff
release. We therefore trained same-recipe models on our own fully documented
pipeline and base all quantitative claims on them, while using the official
weights for high-resolution demonstration.
\paragraph{Quality and trust calibration are separable.}
On BraTS, our self-trained models beat the official weights on PSNR by
$+3$--$4$\,dB, yet the official weights' trust maps localize error better
($+4.80$ vs.\ $+0.36$\,dB sparsification). Better synthesis does not imply
better-calibrated uncertainty: models whose residual error is dominated by
fine texture produce uncertainty signals that are weakly correlated with that
error. Selecting for calibration may hence differ from selecting for quality.
\paragraph{Limitations.}
(1)~The trust map flags \emph{unreliable} regions, not \emph{wrong} ones: a
confidently synthesized hallucination is invisible to it.
(2)~Learnable aggregation overfits 6-subject validation sets; reopening the
learnable route requires $\ge$20 validation subjects.
(3)~The input gate is validated for brain tasks; the pelvis gate is unreliable
and disabled for CT tasks.
(4)~Image-level QC remains unsolved at this scale.
(5)~All outputs are research demonstrations and must not guide clinical
decisions.

\section{Conclusion}
We turned the discarded internals of a state-of-the-art diffusion-bridge
sampler into a pixel-wise trust map, score, and QC workflow for medical image
translation, validated across three datasets and five directions with strong
localization (AUROC up to $0.93$), calibration (separation up to $73\times$),
and clinical reach into tumor patients. Beyond the specific signals, the
broader message is methodological: generative samplers compute more evidence
than they use, and harvesting it converts a synthesis engine into a
trust-assessable system. Future work will reopen learnable aggregation at
larger subject scale, extend the analysis to 256$\times$256 models, and
integrate radiologist feedback.

\paragraph{Acknowledgments.}
We thank the SelfRDB authors (Bilkent ICON Lab) for releasing code and weights
under MIT, and the IXI, BraTS, and SynthRAD consortia for public data. This
work is a student research demonstration; outputs must not be used for
clinical decisions.

\bibliographystyle{unsrtnat}
\bibliography{references}

\end{document}
